\documentclass[11pt,a4paper]{article}

\usepackage[utf8]{inputenc}
\usepackage[T1]{fontenc}
\usepackage{lmodern}
\usepackage[margin=2.5cm]{geometry}
\usepackage{array}
\usepackage{titlesec}
\usepackage[hidelinks]{hyperref}

\newcolumntype{R}[1]{>{\raggedright\arraybackslash}p{#1}}

\titleformat{\section}{\normalfont\large\bfseries}{\thesection}{0.6em}{}
\titlespacing*{\section}{0pt}{1.3ex plus 0.3ex}{0.5ex}

\newenvironment{tightitemize}
  {\begin{itemize}
     \setlength{\topsep}{2pt}\setlength{\itemsep}{0pt}
     \setlength{\parsep}{0pt}\setlength{\parskip}{0pt}}
  {\end{itemize}}

\setlength{\parskip}{0.4em}
\setlength{\parindent}{0pt}

\title{\vspace{-1.5cm}\textbf{Thesis Plan}\\[0.2em]
\large Performance, Safety, and Engineering Trade-offs of Rust and C++\\
in Nokia RAN Control-Plane Systems}

\author{Ha Vien \\
\normalsize University of Helsinki, Master's Programme in Computer Science
(Software study track)\\
\normalsize Nokia instructors: Keijo Heljanko, Christoph Luneburg, Petteri Kotimaa}

\date{\normalsize Revised 6 August 2026}

\begin{document}
\maketitle
\vspace{-0.9em}

\section{Background and motivation}

Modern RAN control-plane software (the L3 of a gNB, e.g.\ the
gNB-CU-CP~\cite{3gpp38401}) is written mostly in C++, which delivers the
throughput and latency such systems need but carries well-known memory-safety
and concurrency risks (use-after-free, data races, double-free). These classes
dominate real defect populations, and migrating to memory-safe languages is now
an explicit policy position~\cite{msrc2019memory,cisa2023memorysafe}. Rust
offers comparable performance with safety guaranteed at compile time on
machine-checked foundations~\cite{jung2018rustbelt}.

Its fit for latency-critical telecom control planes is unproven. The ``Rust
performs like C++'' consensus rests on general-purpose
benchmarks~\cite{pereira2017energy} that report average-case behaviour; closer
studies find real costs, both in networking code~\cite{emmerich2019drivers} and
in specific mechanisms such as bounds and overflow
checking~\cite{zhang2022runtime}. For a RAN control plane the operative metric
is not the average but the tail~\cite{dean2013tail}. This thesis is carried out
at Nokia on a subsystem modelled on a real Nokia gNB-CU-CP component.

\section{Research gap}

Two master's theses at Ericsson are the closest prior work. Karlb\"ack and
Or\"o compare Rust and Go against C++ for RAN applications and conclude C++
remains the choice where performance is critical~\cite{karlback2021ran}.
Arvidsson and Kaller measure median and 99th-percentile latency for 5G
control-plane procedures across five asynchronous frameworks, and find that the
Rust runtimes match C++ at the median but their p99 rises more steeply past a
threshold rate~\cite{arvidsson2025async}.

Neither isolates the language: in both, the runtime and concurrency model vary
together with it, so a measured difference cannot be attributed to the language
rather than to the framework built on it. The gap is therefore a comparison
that holds the specification, algorithms, concurrency model, allocator, and
build settings fixed, \emph{and} sweeps past the saturation knee, \emph{and}
measures what incremental migration costs.

\section{Objective and research questions}

The thesis builds two functionally equivalent implementations of a UE Context
Management service in modern C++ and in Rust and compares them under a
spec-justified synthetic signalling workload. Since similar average throughput
is already known, it targets four less-explored dimensions:

\begin{tightitemize}
  \item[\textbf{RQ1}] How do Rust and C++ compare in latency predictability
  (mean, p99, p99.9, jitter) under increasing load and overload, with the
  concurrency model held fixed?
  \item[\textbf{RQ2}] Which classes of C++ memory and concurrency defects does
  safe Rust eliminate at compile time, at what stage are the rest detected, and
  which classes does it \emph{not} prevent?
  \item[\textbf{RQ3}] What overhead, if any, does Rust's safety model impose on
  throughput, CPU, and memory?
  \item[\textbf{RQ4}] What are the practical implications (effort, FFI
  interoperability, risk) of adopting Rust incrementally in an existing C++ RAN
  codebase?
\end{tightitemize}

RQ2 is deliberately two-sided: a catalogue containing only defects Rust
prevents would confirm its own premise, so it spans the boundary: five
classes safe Rust rejects and two it does not. Memory-safety defects in real
Rust code invariably trace back to \texttt{unsafe}~\cite{xu2021rustcve}, while
deadlock remains an open problem for the language~\cite{zhang2024deadlock}.

\section{Methodology}

A single language-agnostic specification (public API, UE state machine, data
structures, sharding and concurrency model) is frozen before any implementation
code is written, so the comparison measures language rather than design
differences. Both implementations use identical algorithms, a shared-nothing
sharded worker-thread model, the same allocator, and comparable build
optimisation. Holding the concurrency model fixed matters because
thread-per-core structuring itself changes tail
latency~\cite{enberg2019threadpercore}. Both must pass the same parity suite.

The workload is swept from nominal load to overload on pinned,
frequency-locked hardware, with interleaved repeats and bootstrap confidence
intervals~\cite{georges2007rigorous}, latency captured in high-dynamic-range
histograms against an intended send schedule so coordinated omission does not
hide the tail~\cite{tene_hdrhistogram}. Evaluation adds a controlled
fault-injection study~\cite{natella2016fault} and an FFI prototype. On the FFI,
the plan takes a narrow position: confining \texttt{unsafe} to the seam is a
statement about what the compiler checks, not run-time protection, since host
and component share an address space~\cite{rivera2021galeed}. Only synthetic
data is used, so no personal or confidential data is processed.

\section{Expected contributions}

\begin{tightitemize}
  \item A reproducible, telecom-relevant benchmarking framework.
  \item A frozen shared specification and two parity-tested implementations.
  \item A quantitative comparison emphasising tail latency under overload.
  \item A defect-class taxonomy mapping both what the compile-time guarantee
  covers and where it stops.
  \item Evidence-based guidelines for incremental Rust adoption, where the
  dominant cost is the learning curve rather than
  run-time~\cite{fulton2021adopting}.
\end{tightitemize}

\section{Status and timeline}

The implementation and safety milestones landed roughly two months ahead of the
June plan. The single remaining dependency is access to a bare-metal Linux
testbed: the measurements so far ran under WSL2, which exposes neither
frequency governor nor SMT control, so RQ1 and RQ3 are labelled preliminary.
The protocol and campaign driver are written, so this is provisioning work
rather than development.

\begin{center}
\begin{tabular}{@{}R{0.30\linewidth} R{0.63\linewidth}@{}}
\hline\hline
Item & State \\
\hline
Spec + implementations & Done; both pass the shared conformance suite \\
RQ2 (fault injection) & Final: 5 classes compile-time prevented, 2 not \\
RQ4 (FFI prototype) & Final: 12-entry-point seam, core free of \texttt{unsafe} \\
RQ1 / RQ3 (benchmarks) & Harness done; numbers preliminary pending testbed \\
Thesis document & Full draft, builds from the measurement logs \\
\hline
\end{tabular}
\end{center}

\begin{center}
\begin{tabular}{@{}R{0.16\linewidth} R{0.14\linewidth} R{0.63\linewidth}@{}}
\hline\hline
Month (2026) & Goal & What happens \\
\hline
July & Draft & Done: thesis content drafted in full; review and revision pass
over all chapters. \\
August & Measure & Provision the bare-metal testbed; run the full campaign;
regenerate all figures and numbers. \\
September & Resolve & Finalise RQ1/RQ3 analysis; confirm or retract the tail
claim; revise affected chapters. \\
October & Review & Draft to supervisors and Nokia instructors; revision cycle;
artifact packaging. \\
November & Submit & Final revisions, colloquium presentation, submission. \\
\hline
\end{tabular}
\end{center}

\textbf{Risks.} If the testbed is delayed, RQ1 and RQ3 stay preliminary; the
pipeline is fully scripted, so turnaround after access is short, and a cloud
bare-metal instance or a Nokia lab host is the fallback. If the tail difference
proves unresolvable even on bare metal, it is reported as a bounded null result
with the confidence interval stated, which is a legitimate finding. Single-author
skill asymmetry is bounded by the frozen specification and parity suite, and
acknowledged as residual for performance.

\textbf{Out of scope.} The full gNB, PHY/L2 and the user plane, real captured
traces, managed-language alternatives, and production operational concerns.

\bibliographystyle{abbrv}
\bibliography{bibliography}

\end{document}
